\scanpage{31}
Soit $X$ un espace, $(X_i)_{i\geq0}$ une filtration croissante (pour fixer les idées) de $X$ par des sous-espaces fermés. On dit que la \uline{filtration est équisingulière} si $\forall i\geq0$, le diagramme formé de
\diagram{\begin{tikzcd}
(\,X_0 \arrow[r] & X_1 \arrow[r] & \cdots \arrow[r] & X\,)
\end{tikzcd}}
est équisingulier le long de $X_i-X_{i-1}=V_i$.

(NB on pose $X_{-1}=\varnothing$).

Appelons ``feuille (\unsure{entrante})'' de la filtration toute composante connexe d'un $V_i=X_i-X_{i-1}$.
\margin{\unsure{1|?}}
Les feuilles \unsure{entrantes} forment donc une \uline{partition} de $X$ si $X$ est réunion des $X_i$, ce qu'on suppose désormais.

\uline{Proposition} Si la \deleted{\unsure{condition}} filtration est équisingulière, alors l'adhérence de toute feuille est une réunion de feuilles.

(Si $X$ est compact, et les $X_i$ loc. connexes --- ce qui est vrai dans le contexte lin. \unsure{ou} alg. ou an. ``par morceaux'' --- alors il s'agit bien sûr d'une réunion \uline{finie} de feuilles $\ldots$)

Les adhérences des feuilles sont appelées \uline{feuilles fermées}. Il y a corr. biunivoque par $F\longmapsto\overline F$ entre feuilles (\unsure{entrantes} --- on \unsure{ose} les appeler ``ouvertes'') et feuilles fermées. Les parties \uline{fermées} de $X$ qui sont réunions de feuilles sont aussi celles qui sont réunions de feuilles fermées; Toute intersection et toute réunion finie de telles parties est une
\mspage{32}
